\subsection{Task 3: Do learned gate weights matter with fixed experts?}
\label{sec:gate-investigation}
\paragraph{Question and hypothesis.}
Task 2 changed expert identities and weights together. We ask whether the learned
weights still matter when the selected experts are fixed. The hypothesis is that
flattening the weights increases CE and reversing their assignment degrades it
further, indicating sensitivity to which selected expert receives more weight.

\paragraph{Design and controls.}
We use the same validation-selected seed-416 checkpoint without retraining.
Only the last layer's two gate weights change:
\begin{equation}
\widetilde g_e(\lambda)=\tfrac12+\lambda(g_e-\tfrac12),\qquad e\in\{1,2\}.
\end{equation}
The four conditions, specified before evaluation, are original ($\lambda=1$),
uniform ($0$), midpoint ($0.5$), and swapped ($-1$). For two weights summing to
one, swapping gives $(g_2,g_1)$ and preserves their multiset, entropy, and norm.
It changes their assignment to expert identities. Uniformization additionally
changes concentration. All preceding computations, model parameters, ordered
expert IDs, and selected token--expert counts remain fixed. Each condition has
603,264 parameters and 133,120 matrix MACs/token. Thus this is a weighting
intervention rather than a computation-budget comparison.

Evaluation uses the same disjoint test packs: 111,488 text and 1,179,520 code
targets, with batch 32, length 64, and equal domain weighting. The original CE
exactly reproduces Task 1. Hashes match for every last-layer input and ordered
route across conditions; expert counts and checkpoint tensors also match.
Float64 checks verify the weighted-output oracle, 14 executed pairs for seven
tokens, prefix invariance, and independence from other batch rows. Unlike Task 2's
cross-sequence perturbation, these interventions remain token-local.

\begin{table}[ht]
\centering
\small
\caption{Last-layer weight intervention on the fixed seed-416 checkpoint. CE is in nats; $\Delta$ is relative to original mean CE.}
\label{tab:gate-ablation}
\begin{tabular}{@{}lrrrrrr@{}}
\toprule
Weights & $\lambda$ & Text CE & Code CE & Mean CE & $\Delta$ CE & PPL \\
\midrule
Original & 1 & 1.8969 & 1.6289 & 1.7629 & +0.0000 & 5.8292 \\
Uniform & 0 & 2.0792 & 1.8020 & 1.9406 & +0.1777 & 6.9631 \\
Midpoint & 0.5 & 1.9298 & 1.6640 & 1.7969 & +0.0340 & 6.0309 \\
Swapped & -1 & 2.8874 & 2.4569 & 2.6721 & +0.9092 & 14.4706 \\
\bottomrule
\end{tabular}
\end{table}


\paragraph{Results and interpretation.}
Table~\ref{tab:gate-ablation} shows that midpoint, uniform, and swapped weights
increase mean CE by 0.0340, 0.1777, and 0.9092 nats, respectively; both domains
show the same ordering. These measurements support the stated hypothesis at this
checkpoint: preserving expert selection alone does not preserve prediction
quality. Swapping is particularly damaging despite preserving weight concentration,
supporting the importance of matching weights to expert identities rather than
concentration alone. It does not establish that the learned weights are optimal
over all possible assignments.

Figure~\ref{fig:gate-ablation} provides context: original mean gate gaps are
0.8189 for text and 0.7680 for code, with medians 0.9207 and 0.8919. Only 2.20\%
and 3.52\% of tokens have a gap below 0.1. Thus uniformization and swapping
usually make substantial changes, rather than perturbing nearly equal weights.
The observed gap distribution is consistent with the sensitivity, but is not
itself evidence that larger-gap tokens incur larger loss increases.

\begin{figure}[ht]
\centering
\includegraphics[width=\linewidth]{gate_ablation.pdf}
\caption{Fixed-expert weight intervention. (a) Domain CE changes at the four
specified settings; lines connect evaluated points only. (b) Original last-layer
gate-gap distributions, normalized within each domain, with bin width 0.05.
These are single-checkpoint measurements, not averages over training seeds.}
\label{fig:gate-ablation}
\end{figure}
\FloatBarrier

\paragraph{Limitations and reproducibility.}
This experiment concerns one checkpoint and the final layer, with no retraining
or independent training-seed replication. It does not establish the causal
benefit of learning gates during training, nor semantic expert specialization.
The four test evaluations describe pre-specified interventions rather than select
a new deployed model. Run \path{MoE/gate_ablation.py} on CPU; the separate
\path{results/gate_ablation.json} records metrics, checks, and source hashes.
The AI Use Statement below discloses assistance with this investigation.
